\documentclass[11pt,a4paper]{article}
\usepackage[utf8]{inputenc}
\usepackage[T1]{fontenc}
\usepackage{amsmath,amssymb,amsthm}
\usepackage{booktabs}
\usepackage{xcolor}
\usepackage[margin=2.5cm]{geometry}
\usepackage[colorlinks=true,linkcolor=blue!50!black,citecolor=blue!50!black,urlcolor=blue!50!black]{hyperref}

\newtheorem{theorem}{Theorem}
\newtheorem{lemma}{Lemma}
\theoremstyle{remark}
\newtheorem{remark}{Remark}

\newcommand{\R}{\mathrm{R}}
\newcommand{\eqn}{2(x-y)=3z}

\title{The four-colour Rado number of $2(x-y)=3z$ is $226$\\[2pt]
\large A cube-and-conquer computation with checked DRAT proofs}
\author{Takumi Ichikawa\\ \small\texttt{github.com/IchiTaku326}}
\date{22 October 2026\\ \small (not peer reviewed; computation finished 22 October 2026 JST)}

\begin{document}
\maketitle

\begin{abstract}
For a linear equation $E$ and $k\ge 1$, the Rado number $\R_k(E)$ is the least $n$ such that every
$k$-colouring of $\{1,\dots,n\}$ contains a monochromatic solution of $E$.
Chang, De Loera and Wesley \cite{CDW22} computed many such numbers by SAT solving; for the equation
$2(x-y)=3z$ and $k=4$ colours they report only the lower bound $\R_4>225$.
We show that $\R_4(2(x-y)=3z)=226$.
The lower bound is an explicit $4$-colouring of $\{1,\dots,225\}$, checked by an independent program.
The upper bound is a cube-and-conquer proof: $14{,}162$ cubes, each shown unsatisfiable, with every
cube re-solved by CaDiCaL and its DRAT proof accepted by \texttt{drat-trim}.
All programs, logs and the colouring are public.
\end{abstract}

\section{Introduction}
Rado's theorem \cite{Rado33} characterises the linear homogeneous equations that are \emph{regular}, i.e.\ that have a
monochromatic solution in every finite colouring of the positive integers.
For a regular equation $E$ the Rado number $\R_k(E)$ is finite for every $k$.
For the family $a(x-y)=bz$ with $\gcd(a,b)=1$ the paper \cite{CDW22} proves, among other results, $\R_3(a(x-y)=bz)=a^3$ for $a\ge b+2$,
gives the lower bound $\R_k\ge a^k$ for all $k$, and tabulates computed values for $2$, $3$ and (for small $a,b$) $4$ colours.
In their four-colour table the entry for $a=2$, $b=3$ is ``$>225$''; to our knowledge no exact value has been published
(Section~\ref{sec:lit}).
This note closes that entry.

\begin{theorem}\label{thm:main}
$\R_4(2(x-y)=3z)=226$.
\end{theorem}

\paragraph{Contributions.}
The new content is (i) the exact value in Theorem~\ref{thm:main}, (ii) a complete, resumable pipeline whose every step
leaves a time-stamped log, and (iii) a DRAT certificate for each of the $14{,}162$ cases, regenerable in about one hour on a desktop PC.

\paragraph{Conventions.}
Here a \emph{solution} is a triple $(x,y,z)$ of positive integers with $2(x-y)=3z$; the variables need not be distinct
(for example $(4,1,2)$ and $(5,2,2)$ are solutions).
Throughout, ``colouring of $[n]$'' means a map $[n]=\{1,\dots,n\}\to\{0,1,2,3\}$, and it is \emph{good} if it has no monochromatic solution
with $x,y,z\le n$.
Then $\R_4>n$ iff a good colouring of $[n]$ exists, and Theorem~\ref{thm:main} is the conjunction of the two statements below.

\begin{theorem}[Lower bound]\label{thm:lower}
There is a good $4$-colouring of $[225]$.
\end{theorem}

\begin{theorem}[Upper bound]\label{thm:upper}
There is no good $4$-colouring of $[226]$.
\end{theorem}

\section{Lower bound}
A good colouring of $[225]$ was found by a SAT solver (CaDiCaL through PySAT) on 20 October 2026 (JST) and is stored in
\texttt{results/lower\_bound\_n225/}.
It is checked by \texttt{tools/rado\_verify.py}, a program written to be as plain as possible: three nested loops over
$x,y,z\le n$ and a test of $2(x-y)=3z$ and equality of the three colours; it does not share code with the CNF generator.
The check reports \texttt{VALID}.
A good colouring of $[225]$ restricts to a good colouring of every $[m]$, $m\le 225$, so $\R_4>m$ for all $m\le225$.

\section{Upper bound}
\subsection{Encoding}
For $n=226$, $k=4$ we use Boolean variables $v_{j,c}$ ($j\in[n]$, $c\in\{0,1,2,3\}$): ``$j$ has colour $c$''.
The clauses are:
\begin{enumerate}
\item each $j$ gets at least one colour, and at most one colour (6 binary clauses per $j$);
\item for every \emph{solution set} $S=\{x,y,z\}$ and every colour $c$ the clause $\bigvee_{j\in S}\neg v_{j,c}$
 (the set $S$ may have two elements if $z=x$ or $z=y$); there are $8{,}400$ solution sets in $[226]$;
\item symmetry breaking (colours numbered in order of first use): auxiliary variables $u_{j,c}$ meaning ``colour $c$ occurs in $1..j$''
 ($c=0,1,2$) with the defining clauses, together with $v_{1,c}\to\bot$ for $c\ge1$ and $v_{j,c}\to u_{j-1,c-1}$ for $c\ge 1$.
\end{enumerate}
The formula has $1{,}582$ variables and $37{,}891$ clauses.
Every good colouring can be renamed so that colours first appear in the order $0,1,2,3$ (restricted growth); renaming preserves goodness,
and the renamed colouring satisfies the formula. Hence the formula is satisfiable iff a good colouring of $[226]$ exists.

\subsection{Cubes}
Let $d=10$. A \emph{cube} is a restricted-growth string of length $d$ over $\{0,1,2,3\}$ (first letter $0$; each letter at most one larger than the maximum so far) with no
monochromatic solution inside $[d]$. There are exactly $14{,}162$ cubes; this number was obtained by the search program
and re-computed by an independent brute-force enumeration of all $4^{10}$ strings.

\begin{lemma}\label{lem:cover}
If a good colouring of $[226]$ exists, then the formula together with the unit clauses of some cube is satisfiable.
\end{lemma}
\begin{proof}
Rename the colours of a good colouring into restricted-growth form. Its restriction to $[10]$ is a restricted-growth string
without monochromatic solution, i.e.\ one of the cubes, and the colouring satisfies the formula by the previous subsection.
\end{proof}

\subsection{Cube search}
Each of the $14{,}162$ cubes was solved by Glucose (via PySAT, assumptions, conflict budget $2\cdot10^4$, budgets multiplied by $4$ on retry)
in parallel with $11$ workers; all returned UNSAT (21 October 2026, 05:38--09:26 JST, about $41$ CPU-hours in total including retried attempts).

\subsection{Certification}
To avoid trusting an incremental solver, in a second pass (22 October 2026, 01:15--02:23 JST, about $12$ CPU-hours) each cube
was solved again by CaDiCaL as a separate CNF (formula plus unit clauses), producing a DRAT proof, and
\texttt{drat-trim} \cite{WHH14} accepted \emph{every} proof: \texttt{VERIFIED 14162/14162}.
The proofs total about $28$\,GB (largest: $38$\,MB) and are therefore not stored; the second pass regenerates them.
Lemma~\ref{lem:cover} and the $14{,}162$ UNSAT certificates give Theorem~\ref{thm:upper}.

\section{Computational environment}
All computations ran on one desktop PC.
\begin{center}
\begin{tabular}{@{}ll@{}}
\toprule
CPU & Intel Core i5-11400F (6 cores, 12 threads) \\
Memory & 32\,GB DDR4 \\
GPU & NVIDIA GeForce RTX 3060 (not used; CDCL solvers run on the CPU) \\
System & Windows; proof checking in WSL (Linux) \\
Workers & 11 processes for the cube search and for the proof check \\
Lower bound & CaDiCaL 1.5.3 through PySAT \\
Cube search & Glucose 4 through PySAT, assumptions, budget $2\cdot10^4$ conflicts \\
Proof check & CaDiCaL 3.0.1 (static build) and \texttt{drat-trim} (version not recorded) \\
\bottomrule
\end{tabular}
\end{center}

\section{Summary of the computation}
\begin{center}
\begin{tabular}{@{}lr@{}}
\toprule
Variables / clauses of the formula ($n=226$) & $1{,}582$ / $37{,}891$ \\
Solution sets of $2(x-y)=3z$ in $[226]$ & $8{,}400$ \\
Cubes (depth $10$) & $14{,}162$ \\
Cubes UNSAT / SAT & $14{,}162$ / $0$ \\
DRAT proofs accepted by \texttt{drat-trim} & $14{,}162$ / $14{,}162$ \\
Total size of the proofs & about $28$\,GB \\
Cube search, wall clock & 3 h 48 min (11 workers) \\
Proof check, wall clock & about 68 min (11 workers) \\
\bottomrule
\end{tabular}
\end{center}

\section{What this result depends on}
We list the places where an error could still hide.
\begin{enumerate}
\item \textbf{Encoding.} The formula generator is the only place where the problem statement enters the proof. Evidence for its
 correctness: agreement with exhaustive search on small cases; reproduction of published values
 (Schur numbers, $\R_4(x-y=2z)=171$ and the exact $3$-colour values of \cite{CDW22} in the range we tested; $338$ comparisons with
 published values or theorems in a larger table, no disagreement); and validity of the lower-bound colouring under the independent checker.
\item \textbf{Symmetry breaking.} The argument above is a pen-and-paper argument; it was not machine-checked.
\item \textbf{Cube list.} Completeness rests on Lemma~\ref{lem:cover} and the count $14{,}162$.
\item \textbf{Proof checker.} \texttt{drat-trim} is widely used but is itself unverified software; a check with a formally verified
 checker such as \texttt{cake\_lpr} \cite{THM21} has not been done.
\item \textbf{Hardware and software.} Runs were done on a single desktop PC (11 worker processes); the version of \texttt{drat-trim} was not recorded.
\end{enumerate}

\section{Literature status}\label{sec:lit}
We checked, on 22 October 2026, the arXiv listing of \cite{CDW22} (only version v1, 7 October 2022, whose four-colour table
gives $>225$ for $a=2,b=3$), a student report from a SAT group on Rado numbers (three colours only),
and arXiv:2505.12085 (three colours, other equations); none states a value for this equation.
This search was not exhaustive: the list of papers citing \cite{CDW22} and the authors have not yet been consulted.
If the value is already known, this note should be regarded as an independent verification.

\section{Reproducibility}
Code, logs and the colouring: \url{https://github.com/IchiTaku326/rado-226} (MIT licence); archived copy of release v1.0: \url{https://doi.org/10.5281/zenodo.23565748}.
Step 0 validates the colouring (seconds), step 1 reruns the cube search (hours on about 11 cores, resumable),
step 2 reruns the DRAT check (about one hour on 11 cores).
Every row of \texttt{results/} carries the JST time at which it was written.

\section*{Acknowledgements and disclosure}
The programs, the computation plan and parts of the text were developed with the help of an AI assistant (Claude, Anthropic).
The author is responsible for the content. This is a small computational result; it has no bearing on the P versus NP problem.

\begin{thebibliography}{99}
\bibitem{CDW22} Y.~Chang, J.\,A.~De Loera, W.\,J.~Wesley, \emph{Rado numbers and SAT computations}, arXiv:2210.03262 (2022), DOI 10.1145/3476446.3535494.
\bibitem{Rado33} R.~Rado, Studien zur Kombinatorik, \emph{Math.\ Z.} \textbf{36} (1933), 424--470.
\bibitem{HKM16} M.\,J.\,H.~Heule, O.~Kullmann, V.\,W.~Marek, Solving and verifying the Boolean Pythagorean triples problem via cube-and-conquer, \emph{SAT 2016}, LNCS 9710, 228--245.
\bibitem{Heule18} M.\,J.\,H.~Heule, Schur number five, \emph{Proc.\ AAAI 2018}; arXiv:1711.08076.
\bibitem{WHH14} N.~Wetzler, M.\,J.\,H.~Heule, W.\,A.~Hunt Jr., DRAT-trim: efficient checking and trimming using expressive clausal proofs, \emph{SAT 2014}, LNCS 8561, 422--429.
\bibitem{THM21} Y.\,K.~Tan, M.\,J.\,H.~Heule, M.~Myreen, cake\_lpr: verified propagation redundancy checking in CakeML, \emph{TACAS 2021}, LNCS 12652.
\bibitem{Biere20} A.~Biere, K.~Fazekas, M.~Fleury, M.~Heisinger, CaDiCaL, Kissat, Paracooba, Plingeling and Treengeling entering the SAT Competition 2020, \emph{SAT Competition 2020: Solver and Benchmark Descriptions}, 51--53.
\bibitem{AS09} G.~Audemard, L.~Simon, Predicting learnt clauses quality in modern SAT solvers, \emph{IJCAI 2009}, 399--404.
\bibitem{IMM18} A.~Ignatiev, A.~Morgado, J.~Marques-Silva, PySAT: a Python toolkit for prototyping with SAT oracles, \emph{SAT 2018}, LNCS 10929, 428--437.
\bibitem{LR14} B.~Landman, A.~Robertson, \emph{Ramsey Theory on the Integers}, 2nd ed., AMS, 2014.
\end{thebibliography}
\end{document}
